\section{The double Gaussian operator of type B }
In this non-commutative setting, random variables are understood to be the elements of the
$\ast$-algebra generated by $\{\B(x\otimes y) ,\B^\ast(x\otimes y)\mid x\otimes y \in \HH_\R\}$. Particularly interesting are their mixed moments.
In order to work effectively on this object we need to define the corresponding statistics. We provide an explicit formula for the combinatorial moments, involving the number of crossings and negative pairs of a partition. First, we need to define the operators, the set of partitions and statistics of type B.

\subsection{The orthogonal polynomials}
\begin{Definition}
	The operator
	\begin{equation*}
		\G(x\otimes y)= \B(x\otimes y) +\B^\ast(x\otimes y),\qquad x\otimes y \in \HH_\R
	\end{equation*}
	on $\F$ is called the \emph{double Gaussian operator of type B}.
Denote by $\state$ the vacuum vector state $\state(\cdot)=\langle\Omega \otimes \Omega, \cdot\text{ } \Omega\otimes \Omega\rangle$.
\end{Definition}
